\item \textbf{Ilusión óptica del mirador de la fresa (Mirage).}
Se logra un efecto tridimensional sorprendente de una fresa flotante usando dos espejos parabólicos cóncavos idénticos orientados cara a cara, cada uno de distancia focal $7.50\text{ cm}$, colocados con sus vértices separados $7.50\text{ cm}$. Al colocar una fresa real en el fondo del espejo inferior, se forma una imagen de ella en la abertura superior en el centro del espejo de arriba.
\begin{enumerate}
    \item Demuestre que la imagen final se forma efectivamente en dicha posición y describa sus características físicas (real/virtual, vertical/invertida, magnificación).
    \item Un efecto impresionante ocurre al apuntar una linterna a la imagen tridimensional desde un ángulo inclinado: ¡parece que el haz se refleja en la fresa inexistente! Explique físicamente este fenómeno.
\end{enumerate}